%------------------------------------------------------------------------------%
\section{Integrais Definidas}
\label{sec:Integrais_Definidas}

Uma operação relacionada é o cálculo da área entre o gráfico da função $f$ e
o eixo-$x$, chamamos essa operação de integral definida como descrito na
próxima definição.

%------------------------------------------------------------------------------%
\begin{definicao}{Integral Definida}{def:integral-defnida}
  A \emph{Integral Definida}\index{Integral!definida} de uma função real $f$
  é a área, considerando o sinal, entre o gráfico da função e o eixo-$x$
  no intervalo $[a, b]$
  \[
  \int_{a}^{b} f(x) \, dx
  \]
  Essa área é o limite da \emph{Soma de Riemman} correspondente.
\end{definicao}
%------------------------------------------------------------------------------%

A integral definida possui as propriedades descrita na proposição a seguir.

%------------------------------------------------------------------------------%
\begin{proposicao}{Propriedades da Integral Definida}{pro:propriedade-integral-definida}
  Dadas funções integráveis $f$ e $g$ e a constante $c$ temos que
  \begin{enumerate}
    \item \( \int_a^b c\,dx = c(b-a) \)
    \item \( \int_a^b f(x)+g(x)\,dx = \int_a^b f(x)\,dx + \int_a^b g(x)\,dx\)
    \item \( \int_a^b f(x)-g(x)\,dx = \int_a^b f(x)\,dx - \int_a^b g(x)\,dx\)
    \item \( \int_a^b cf(x)\,dx = c \int_a^b f(x)\,dx \)
    \item \( \int_a^c f(x)\,dx + \int_c^b f(x)\,dx = \int_a^b f(x)\,dx \)
  \end{enumerate}
\end{proposicao}
%------------------------------------------------------------------------------%

\hrulefill

Interpretação Geométrica da integral definida.
Considere $f$ uma função não negativa no intervalo $[a,b]$ conforme a
Figura~\ref{fig:integral_definida_1}.

\begin{figure}
  \centering
  \includegraphics[width=6cm,height=5cm]{integral_definida_01}
  \caption{Gráfico da função $y=x^2$.}
  \label{fig:integral_definida_1}
\end{figure}

Então a integral $\int_a^b f(x)dx$ significa a área sob a curva $y=f(x)$
no intervalo $[a,b]$, conforme indica a Figura~\ref{fig:integral_definida_2}.

\begin{figure}
  \centering
  \includegraphics[width=6cm,height=5cm]{integral_definida_02}
  \caption{Gráfico da função $y=x^2$.}
  \label{fig:integral_definida_2}
\end{figure}

Se $f(x)$ é uma função que muda de sinal no intervalo $[a,b]$,
então $\int_a^b f(x)dx$ pode ser interpretada como a diferença
de áreas conforme ilustra a Figura~\ref{fig:integral_definida_3}.

\begin{figure}
  \centering
  \includegraphics[width=6cm,height=5cm]{integral_definida_03}
  \caption{Gráfico da função $y=x^2$.}
  \label{fig:integral_definida_3}
\end{figure}

%------------------------------------------------------------------------------%
\begin{teorema}{Propriedades da integral definida}{teo:propriedades-integral-definida}
  \begin{enumerate}
    \item \(\int_{a}^b f(x)dx= -\int_{b}^a f(x)dx\)
    \item \(\int_{a}^a f(x)dx=0\)
    \item \(\int_{a}^b cdx = c(b-a)\), $c$ é uma constante.
    \item \(\int_{a}^b [f(x)+g(x)]dx = \int_{a}^b f(x)dx + \int_{a}^b g(x)dx\)
    \item \(\int_{a}^b cf(x)dx= c \int_{a}^b f(x)dx\), $c$ é uma constante.
    \item \(\int_{a}^b [f(x)-g(x)]dx= \int_{a}^b f(x)dx- \int_{a}^b g(x)dx\)
    \item \(\int_{a}^b f(x)dx=\int_{a}^c f(x)dx+ \int_{c}^b f(x)dx\)
    \item Se $f(x)\geq 0$, $a \leq x \leq b$, então \(\int_{a}^b f(x)dx \geq 0\).
    \item Se $f(x) \geq g(x)$, $a \leq x \leq b$, então \(\int_{a}^b f(x)dx \geq \int_{a}^b g(x)dx\)
    \item Se $m \leq f(x) \leq M$, para $a \leq x \leq b$ então
    \[
    m(b-a) \leq \int_{a}^b f(x)dx \leq M(b-a)
    \]
  \end{enumerate}
\end{teorema}

%------------------------------------------------------------------------------%
\begin{example}
  Calcule a integral da função de Heaviside
  no intervalo $[-1, 2]$.
  \solution
  Primeiro precisamos relembrar a definição da função de
  Heaviside~\eqref{eq:funcao_heaviside}
  \[
    \heaviside(x) = \begin{cases}
      \; 0 & \; x < 0 \\
      \; 1 & \; x > 0
    \end{cases}
  \]
  Notamos que essa função é definida por partes, assim precisamos usar
  a propriedade 7 do Teorema~\ref{teo:propriedades-integral-definida}
  para dividir a integral em duas partes
  \[
    \int_{-1}^{2} \heaviside(x) \,dx
    = \int_{-1}^{0} 0 \,dx
    + \int_{0}^{2} 1 \,dx \,.
  \]
  Agora precisamos calcular separadamente as áreas entre cada função e o
  eixo-$x$.

  Na primeira integral a função é sempre zero e portanto a área também
  será zero
  \[
    \int_{-1}^{0} 0 \,dx = 0 \,.
  \]
  Na segunda integral a função é constante igual a 1, portanto a área entre
  o gráfico e o eixo será a área do retângulo de altura 1 e base igual
  a o comprimento do intervalo de integração
  \begin{align*}
    \int_{0}^{2} 1 \,dx
    & = \int_{0}^{2} dx  \\[3mm]
    & = \text{base} \times \text{altura} \\[3mm]
    & = (2-0)\times 1 = 2
  \end{align*}
  Concluímos que
  \[
    \int_{-1}^{2} \heaviside(x) \,dx = 0 + 2 = 2 \,.
  \]
\end{example}
%------------------------------------------------------------------------------%


%------------------------------------------------------------------------------%
